% Шаблон плана: «План управления конфигурацией ПО», «План разработки ПО», «План верификации ПО».
% Скопируйте в docs/0N_<имя>.tex, замените название, версию и дату, заполните разделы.
% Красные пометки \clarify{…} — места, которые надо заполнить; перед сдачей их не остаётся.
% Правила оформления — docs/00_formatting.tex. Собрать: ./build.sh docs/0N_<имя>.tex
\documentclass{tpks}
\doctitle{План …}
\docversion{1}
\docdate{ДД.ММ.ГГГГ}

\begin{document}
\makeheader
% Лист регистрации изменений — сразу после \makeheader. Новая версия — новая строка.
\begin{changelog}
\change{1}{ДД.ММ.ГГГГ}{Первая версия}{Фамилия И. О.}{Задание к паре ДД.ММ.ГГГГ}
\end{changelog}

\section{Введение}
\subsection{Назначение}
\clarify{цель плана}

\subsection{Область применения}
\clarify{к каким процессам и данным жизненного цикла относится план}

\subsection{Термины и сокращения}
Определения взяты из общего глоссария проекта.
\printterms{ЕК, Базовая версия}

\subsection{Ссылки}
\begin{enumerate}
  \item \clarify{«Название документа». Что из него используется}
\end{enumerate}

\section{Организация и роли}
\begin{tabl}{|P{4.5cm}|L|P{4cm}|}{Роли участников}[tab:roles]{\textbf{Роль} & \textbf{Обязанности} & \textbf{Участник}}
\clarify{роль} & \clarify{обязанности} & \clarify{ФИО} \\ \hline
\end{tabl}

\section{Процессы}
% Каждый процесс описывается по одной схеме. Входы и выходы — README.md §3.1.
\subsection{\clarify{Название процесса}}
\subsubsection{Цели}
\clarify{}
\subsubsection{Входные данные}
\clarify{}
\subsubsection{Условия входа}
\clarify{}
\subsubsection{Описание}
\clarify{}
\subsubsection{Условия выхода}
\clarify{}
\subsubsection{Выходные данные}
\clarify{}
\subsubsection{Роли}
\clarify{}

\section{Среда и инструменты}
\clarify{средства, которые используются в процессах плана}

\end{document}
